\subsection{How can we Hear Sverdrup and Evjen}

We cannot understand Sverdrup by reading Sverdrup, as we no longer have the common frame of reference. As evidence, consider this clause from P1, "Kingdom of God on earth." I ask you to take it seriously with all the implications at both the level of congregation and sovereign individual operating within that congregation. 

This is a good time to stop and find an old hymnal. Immerse yourself in the Kingdom described in these hymns:

\begin{itemize}
    \item \emph{Battle Hymn of the Republic}
	\item \emph{Stand Up, Stand Up for Jesus}
	\item \emph{We’ve a Story to Tell to the Nations}
	\item \emph{Onward, Christian Soldiers}
	\item \emph{Bringing in the Sheaves}
	\item \emph{From Greenland’s Icy Mountains}
    \item \emph{Work, for the Night Is Coming}
\end{itemize}

There is an incredible optimism woven into these hymns, an American heat whose supreme confidence borders on sinful pride, asserting that the Kingdom will come to all the earth. Do we dare take the next step and declare that the "Kingdom of God on earth" is here embodied in the congregation and its faithful members. Is this not the natural natural force of P1? 


\begin{quote}
According to the Word of God, the congregation is the right form of the Kingdom of God on earth."
\end{quote}

You can read more about this uniquely American church history era under the category of postmillennialism. You will find the power exemplified in the refrain of "We’ve a Story to Tell to the Nations":

\begin{quote}
For the darkness shall turn to dawning,
and the dawning to noonday bright,
and Christ's great kingdom shall come on earth—
the kingdom of love and light.
\end{quote}

This is the providential calling to the congregation as well as individual duty, urgency, and frankly, the unquestioned march driving the church's mission forward. Understand that this is the empowering water that surrounds Sverdrup, Oftedahl, Evjan, and the Congregationalists in the mid-to-late 1800s. Like a fish, I'm not convinced they could see it. Yet, it is there with all the twists and turns that carry "John Brown's Body" to your hymnal with the words: 

\begin{quote}
Glory! Glory! Hallelujah!
Glory! Glory! Hallelujah!
Glory! Glory! Hallelujah!
His truth is marching on.
\end{quote}

Be careful as I'm not saying Sverdrup and Oftedahl should be written off as postmillennialism thinkers. To the contrary, I believe they offer something profound that calls to us today. 


\subsubsection{A model of the Congregation}

To better understand this world I will offer the model of the visible congregation as a ship on a mission at sea. Keep in mind that the ship is my metaphor and, to the best of my knowledge, not used by either Sverdrup or Evjan. 

We must be able to see that water in which the visible congregation is sailing. 


manning a ship.

Evjen was writing about Sverdrup in the decade following his untimely death. We are reading Evjen and Sverdrup a century later. I suspect Evjen was too close to the fire and we are so far away that we cannot distinguish the smoke through the fog. 

In the previous section I was very careful to define terms of dogmatics, ecclesiology, and polity. We also smuggled in a few Sverdrup terms including principle 1 and Christ's free gift of spirit and life.



This is not Sverdrup or Evjen language. It's not doctrine. It's a model. And like all models is useful but imperfect. 

However, the model provides a conceptual framework that make it easier to visualize the congregation. Conversations are easier, and operational failures are easier to describe and impossible to deny.
  
I was a sailor who sailed on icebreakers. I've been far north of Alaska and south to Antarctica. There is an only saying that a "ship is safe in port, but that's not what a ship is for."

That is important as we are going to describe the congregation the congregation as ship on a mission with a calling from Christ. Instead of using "congregation as the body of Christ with each of us a members," we describe members of the congregation as member of the ship's crew. In this case, the ship is the visible structure in which Christ's work is carried out. 

This is not a safe ship. It's not some sentimental ark to weather the storm. It's an operation ship dependent on the shared responsibility of the crew to perform the ship business while taking care of all members.

This is not a healthy ship. We are all sinners, Satan is whispering to the crew, and the real world presses in with political, societal, and economic constraints. There are also self professed enemies to Christianity. 

A congregation is safe when it avoids mission, repair, conflict, responsibility, and pays the pastor until it closes its doors. But that's not what a congregation is for. Or as Sverdrup once stated, a congregation may appear to have all of it's members, but it may be missing the most important things, and that is life.  



In my opinion, P1 is not merely an academic problem as framed by Evjen, but an existential vocational problem. At its core it demands something from every member of the crew. If they are unwilling or unable to develop and use their gifts the congregation fails. "Spirit and life" is not a revival statement about "Jesus Loves me," but a life and death question about the duty performed by the individual members of the congregation. In biblical terms, P1 is a declaration of works follow faith. A serious acceptance of "Kingdom of God on earth" demands no less.  

And finally we arrive at Sverdrup's core message concerning, a living and visible congregation acts with all hands active.

As Evjen stated, "all gifts and use for all hands." FACTCHECk 

Can we bear the weight of our earthly vocation? Will we accept P1 not as a sentimental statement of spirit and life statement or a serious call to act as serious member of the body of Christ? 

Can P1 exist as a living confession without the fruits of "every single man and women who is called in Christ's name"? 


And so we introduce Evjen with his own words:

\begin{quote}
The question—the real one—becomes: Does one stand in living accord with the future of God’s kingdom, and is one willing to let oneself be used for the kingdom’s advance, to labor for the fulfillment of the prayer: hallowed be your name; your kingdom come; as in heaven so also on earth—your will be done!”
\end{quote}  

In the Reformation tradition we acknowledge that faith alone justifies, but faith is never alone.

\begin{quote}
“My language has been the bracing language of duty.”
\end{quote} 



\begin{quote}
“It is possible to bury the principles and get a peace that is as silent as the grave’s.”
\end{quote} 





